% Generated draft: requires mathematical and source review.
\documentclass[11pt]{article}
\usepackage[margin=25mm]{geometry}
\usepackage{fontspec}
\setmainfont{Latin Modern Roman}
\usepackage{amsmath,amssymb,mathtools,unicode-math}
\setmathfont{Latin Modern Math}
\usepackage{enumitem,needspace,xurl,hyperref}
\setlength{\parindent}{0pt}
\setlength{\parskip}{5pt}
\setlength{\emergencystretch}{3em}
\newcommand{\N}{\mathbb{N}}
\newcommand{\Z}{\mathbb{Z}}
\newcommand{\Q}{\mathbb{Q}}
\newcommand{\R}{\mathbb{R}}
\newcommand{\C}{\mathbb{C}}
\newcommand{\F}{\mathbb{F}}
\newcommand{\A}{\mathbb{A}}
\newcommand{\PP}{\mathbb{P}}
\newcommand{\E}{\mathbb{E}}
\newcommand{\eps}{\varepsilon}
\newcommand{\dd}{\,\mathrm{d}}
\begin{document}
\section*{UnsolvedMath ID: 30002321}
\textbf{Source problem code:} catalogue entry
\subsection{Definitions and mathematical statement}
Under the conventional real interpretation, a nonnegative matrix is a matrix \(A\in\R_{\geq 0}^{n\times n}\). Its ordinary rank \(\operatorname{rank}_{\R}(A)\) is its rank over \(\R\), and its nonnegative rank is
\[
\operatorname{rank}_{+}(A)=\min\left\{r\in\Z_{\geq 0}: A=UV,\quad U\in\R_{\geq 0}^{n\times r},\quad V\in\R_{\geq 0}^{r\times n}\right\}.
\]
The case \(r=0\) represents the zero matrix. The supplied statement does not explicitly specify the scalar field; the definitions here use real matrices and real nonnegative factorizations.

Interpreting the asymptotic question as asking for examples in arbitrarily large dimensions, does there exist an integer \(R\geq 1\), a constant \(c>0\), and an infinite set \(S\subseteq\{1,2,\ldots\}\) such that, for every \(n\in S\), there exists \(A_n\in\R_{\geq0}^{n\times n}\) satisfying
\[
\operatorname{rank}_{\R}(A_n)\leq R,
\qquad
\operatorname{rank}_{+}(A_n)\geq cn?
\]
Here \(R\) and \(c\) are independent of \(n\). This makes precise one interpretation of bounded ordinary rank and nonnegative rank \(\Omega(n)\). The source does not explicitly distinguish this interpretation from the stronger requirement that such matrices exist for every sufficiently large integer \(n\).

The source also asks, in particular, whether rank-three matrices can have nonnegative rank equal to their dimension. For a specified integer \(n\geq3\), the corresponding condition is
\[
\exists A\in\R_{\geq0}^{n\times n}:\quad
\operatorname{rank}_{\R}(A)=3,
\qquad
\operatorname{rank}_{+}(A)=n.
\]
In the asymptotic context, the intended family question is whether this condition holds for arbitrarily large \(n\); whether the source requires every sufficiently large \(n\) is not explicit. Merely asking for one dimension would admit \(n=3\) and would not capture the intended asymptotic issue.

The supplied dated curation labels the problem partially solved and reports that Shitov ruled out the rank-three, nonnegative-rank-equals-\(n\) family for sufficiently large \(n\). It also states that this result does not settle the broader bounded-rank linear-growth question. These are supplied literature qualifications, not independently verified conclusions of this draft.
\subsection{Short English statement}
Can nonnegative square matrices have ordinary rank bounded by a fixed constant while their nonnegative rank grows at least linearly with their dimension? The source singles out rank-three matrices with nonnegative rank exactly \(n\). Its supplied dated assessment reports an obstruction to that special family in sufficiently large dimensions, while leaving the broader question unsettled by that result. The scalar field and the precise quantification over dimensions require confirmation.
\subsection{Sources}
\begin{itemize}
\item Supplied local UnsolvedMath registry record: permanent ID \(30002321\), title \emph{Large Nonnegative Rank at Bounded Ordinary Rank}, source problem code \texttt{catalogue entry}, in the Oberwolfach Reports Open Problems collection. The record lists proposed year \(2013\) and names no proposer.
\item Supplied report citation: \emph{Combinatorics and Probability (2014)}, Oberwolfach Reports, pp.\ 1087--1152, DOI \texttt{10.4171/owr/2013/18}. \url{https://doi.org/10.4171/owr/2013/18}. The supplied background describes these reports as documenting mathematical workshops and their open-problem discussions. The citation's \(2014\) designation and its \(2013\) DOI component are retained without reconciliation.
\item Supplied literature reference: Shitov, ``upper bound for nonnegative rank.'' \url{https://doi.org/10.1016/j.jcta.2013.10.004}. No fuller bibliographic details are supplied.
\item Supplied literature reference: ``Oberwolfach combinatorics report.'' \url{https://doi.org/10.4171/owr/2013/18}.
\item Literature assessment checked at \texttt{2026-08-21}: the supplied curation labels the problem partially solved, reports that Shitov disproved the proposed rank-three family for sufficiently large \(n\), and says the broader bounded-rank linear-growth question is not settled by that result. This is supplied dated curation, not independently verified literature assessment; the record itself requests independent verification of time-sensitive status and cited claims.
\item The record labels the statement \texttt{corrected\_verified} and reports that its mathematical content was checked against the official Oberwolfach report, with notation, context, and question form normalized. This is a supplied verification claim; no source was independently consulted for this draft.
\end{itemize}
\end{document}
